\section{Notas de Revisión Litográficas (Observaciones del Profesor)}

De acuerdo con las anotaciones en lápiz y marcas de revisión que el profesor José Romero plasmó sobre el manuscrito original, se extraen los siguientes puntos a tomar en cuenta en el documento digital:

\begin{enumerate}
    \item \textbf{Puntuación por Secciones:} El primer bloque (modelado y análisis clásico) cuenta con una anotación de $6/10$, mientras que el desarrollo por variables de estado ostenta un $8/10$. La entrega cierra con la marca aprobatoria del docente.
    
    \item \textbf{Base de Tiempo de Muestreo:} En los cálculos y gráficas de control digital, ratificar que el período de muestreo está fijado exactamente en $T_s = 0.01 \, \text{seg}$ (frecuencia de cálculo de la rutina a $1000 \, \text{Hz}$) para evitar desfases en las ecuaciones de diferencias de los observadores.
    
    \item \textbf{Tratamiento Dinámico de la Barra:} Se destaca el recordatorio de tratar el sistema bajo la perspectiva rigurosa de un \textbf{péndulo físico} en lugar de una masa puntual (péndulo simple) al deducir el momento de inercia ($I = \frac{1}{3}ml^2$) y el centro geométrico de la masa, asegurando la concordancia entre los autovalores teóricos y el comportamiento real del prototipo.
    
    \item \textbf{Aclaratoria en Gráfica de Validación:} El profesor identificó una discrepancia visual en los resultados temporales adjuntos, aclarando que el periodo real de la oscilación amortiguada del sistema físico es de \textbf{16 segundos}, corrigiendo la estimación inicial de 10 segundos.
\end{enumerate}
